\subsection{Model Comparison Framework}

The methodology uses a common physical input-output structure to compare several predictive model families. The comparison is not intended to prove that one model fully solves the thermoelastic wave propagation problem. Instead, the models are treated as alternative predictive components that operate within the same physically motivated framework. This makes it possible to compare their behavior using consistent material parameters, thermal conditions, source--probe geometry, and output interpretation.

The first model family is the Physics-Informed Neural Network, or PINN. In this thesis, the PINN is treated as a physics-informed baseline because it is designed to incorporate physical relations into the learning process. Its role in the comparison is not to claim a complete solution of the full three-dimensional coupled thermoelastic system, but to provide a reference model whose structure is motivated by governing equations and physical consistency.

The second model family is the Fourier Neural Operator, or FNO. Operator-learning methods are relevant because thermoelastic wave propagation can be interpreted as a mapping from input fields, material parameters, and conditions to output response fields or response quantities. In this thesis, the FNO component is considered as an alternative predictive approach within the same input-output contract. The detailed implementation and service structure are discussed later, while this chapter focuses only on the methodological role of the model.

The third model family is MeshGraphNet. Graph-based models are useful for problems where geometry, local neighborhoods, or mesh-like relationships are important. In the context of geological media, this is relevant because spatial structure and directional relationships influence wave behavior. Here, MeshGraphNet is treated as another candidate predictive component rather than as a complete geological simulator.

The fourth component is a Transformer-style baseline. This model family is included as an experimental comparative baseline for structured input-output prediction. Its role is to provide an additional point of comparison against the more physics-oriented or geometry-oriented approaches. As with the other model families, it is evaluated within the same physically motivated input-output framework.

These approaches are compared as alternative predictive components within the same physically motivated input-output framework. The main purpose of the comparison is to examine how different model families respond to the same material parameters, thermal conditions, and propagation directions. Detailed software architecture, API behavior, service-level implementation, and deployment structure are not part of this chapter and are reserved for the implementation chapter.
